\section{Discussion}\label{sec:discussion}

\hw{this section argues from the simulation and the literature only;
rewrite from the measurements}

\textbf{What the simulation can and cannot say.} In the simulation the
prior's value is concentrated below 32 segments, and a one-robot prior
is within 30 \% of a four-robot prior from eight segments on, because
the population was drawn from the model family. On real robots the risk
is misspecification: a robot whose dead zone or reversal point falls
outside the family would be pulled towards the prior and look calibrated
when it is not. The predictive-interval coverage and the held-out
measured segments detect this; the full grids of phase B locate it.

\textbf{Drift.} The day component of the decomposition, the battery term
of the \texttt{bat} variant and the three-day repeat on robot \#1 are
designed to quantify what the Kilobot literature reports qualitatively
\cite{raoufi2023individuality, rubenstein2014kilobot}. If drift turns
out to be a sizeable fraction of the robot component, a recursive
least-squares update of the thirteen parameters during operation,
keyed by chip serial, is one extension
\cite{martinelli2007selfcalibration, roy1999selfcalibration}.

\textbf{Outlook: setting heterogeneity.} If the measured robot
component dominates the decomposition, calibration could remove it, and
the same machinery could set a chosen heterogeneity: Giomi et al.
\cite{giomi2013bristlebots} show that a per-robot spin-to-walk ratio
selects the collective state, and Siebers et al.
\cite{siebers2023disorder} exploit disorder deliberately. Whether a
calibrated GRASPion fleet can serve as an active-matter experiment with
a controlled disorder parameter depends on the measured decomposition.

\textbf{Open questions the data will decide.} Whether the two motors add
linearly, as Notomista et al. \cite{notomista2019brushbot} derive, or
interact; whether transient data identify the CW / CCW split that
steady-state data cannot; whether the PWM frequency changes the map's
shape or only its gain; and whether the onboard magnetometer, once its
hard-iron offset \cite{vasconcelos2011magnetometer} and the offset
\(b(L,R)\) that the running motors induce at each drive level are
mapped, can carry the yaw-rate part
of recalibration without a camera.
